\documentclass[11pt]{article}

\usepackage[utf8]{inputenc}
\usepackage[T1]{fontenc}
\usepackage{lmodern}
\usepackage{geometry}
\usepackage{amsmath,amssymb,amsfonts,amsthm,mathtools}
\usepackage{bm}
\usepackage{enumitem}
\usepackage{microtype}
\usepackage{hyperref}
\usepackage{titlesec}
\usepackage{booktabs}
\usepackage{array}
\usepackage{setspace}
\usepackage{mathrsfs}
\usepackage{physics}

\geometry{margin=1in}
\hypersetup{colorlinks=true, linkcolor=black, urlcolor=blue, citecolor=black}
\setlist[enumerate]{topsep=0.4em,itemsep=0.35em}
\setlist[itemize]{topsep=0.3em,itemsep=0.25em}
\titleformat{\section}{\large\bfseries}{\thesection.}{0.5em}{}
\titleformat{\subsection}{\normalsize\bfseries}{\thesubsection.}{0.5em}{}
\setstretch{1.06}

\newcommand{\Aether}{\AE{}ther}
\newcommand{\Aflow}{\AE{}ther-flow}
\newcommand{\TheoryName}{\Aflow{} Interpretation of Relativity}
\newcommand{\TheoryProgram}{\Aether{} / \Aflow{} framework}
\newcommand{\CompletedMetric}{g^{\sharp}}

\newcommand{\ProjectionOperatorCore}{\mathfrak S^{\mathrm{disp,resp,projop}}}
\newcommand{\CurrentBodyImageCore}{\mathfrak S^{\mathrm{disp,resp,cbimg}}}
\newcommand{\PullbackBodyResponseCore}{\mathfrak S^{\mathrm{sub,proto,pppresec,pb,bodyresp}}}

\newtheorem{definition}{Definition}
\newtheorem{proposition}{Proposition}
\newtheorem{remark}{Remark}
\newtheorem{corollary}{Corollary}
\newtheorem{theorem}{Theorem}

\title{\textbf{The \TheoryName{}}\\[0.4em]
\large Primitive-Reservoir Observer-Reduction\\
\large Section-Centered Hidden-Response\\
\large Canonical Current-Body Image Core\\
\large Next Benchmark Criterion}
\author{}
\date{}

\begin{document}

\maketitle

\begin{abstract}
The current active-primary theorem now proves that the full section-centered hidden-response projection-operator core still carries benchmark-facing presentational redundancy. Once the already-fixed current displacement body, the surjective displacement projection, the coercive positive self-adjoint substrate-response operator, and benchmark faithfulness are fixed, the full listed substrate body is no longer independently live at current scope: all current benchmark-facing uses factor through the canonical image of the fixed current body under the forced minimal lift, and off-image kernel directions are benchmark neutral. The live unexplained stopping point on the recorded route is therefore no longer the full projection-operator core \(\ProjectionOperatorCore\). It is the smaller canonical current-body image core \(\CurrentBodyImageCore\).

This note records the routing consequence of that gain. The next honest benchmark-facing manuscript must now derive or force \(\CurrentBodyImageCore\) itself, or some nontrivial constituent or relation inside it, from still more primitive explicit substrate structure in a way that changes benchmark-facing status. The claim boundary remains fixed: the one operative metric is still exactly \(\CompletedMetric\), the ordinary matter route is unchanged, there is no second operative metric, the low-energy matter law is not enlarged, and no stronger promotion language is licensed. The positive result is therefore routing discipline below the new current-body-image-core frontier.
\end{abstract}

\section{Introduction}

The active derivational program of \TheoryName{} has again moved the live burden upstream by an honest smaller theorem. The previous current theorem showed that the full displacement-response substrate law collapses to the smaller current-body-realizing projection-operator core \(\ProjectionOperatorCore\). The new theorem sharpens that result strictly inside the core itself: the full substrate-body presentation is no longer primitive either. What now remains live at current scope is only the smaller canonical current-body image core \(\CurrentBodyImageCore\).

That changes what now counts as honest primary progress. Another manuscript that merely carries the same canonical image core together with a restored full substrate body, or that re-expands back to the full projection-operator core or the full displacement-response substrate law as if the off-image body directions were still independently live, no longer removes any remaining benchmark-facing freedom. The next honest move must therefore target the smaller image core itself or some genuinely smaller theorem inside it.

\paragraph{Framework and claim boundary.}
\TheoryProgram{} denotes the broader \Aether{} / \Aflow{} framework built on the claims that the \Aether{} is the underlying four-dimensional substrate of reality, the \Aflow{} is its intrinsic ordered motion, observed three-dimensional space is the local experiential slice of that deeper substrate, S-time is the experienced order of change arising from matter, light, and the \Aflow{}, and observed expansion is the three-dimensional appearance of deeper four-dimensional ordered motion. Within that framework, \TheoryName{} denotes the exact-closure theory in which the effective gravitational dynamics are adopted to be exactly Einsteinian with universal matter coupling. In the active sequence, \emph{adoption} means use of that established relativistic dynamics without claiming substrate derivation, while \emph{derivation} is reserved for a first-principles recovery from explicit substrate variables.

The present note stays strictly inside that boundary. It records only which kind of next manuscript now counts as an honest benchmark-facing gain below the new current-body-image-core frontier.

\section{Fixed Inputs}

\begin{definition}[Current-body-image-core frontier inputs]
The criterion recorded here uses the following fixed inputs.
\begin{enumerate}
    \item The benchmark exact-closure package, which fixes the public one-metric GR target of the repository.
    \item The benchmark gatekeeping layer, including the recorded safeguard that blocks same-output deeper-origin relays from counting as new front-line progress once the live benchmark-relevant datum is unchanged.
    \item The earlier theorem proving that the current observer-relevant pullback body-response core \(\PullbackBodyResponseCore\) is forced by a benchmark-faithful section-centered displacement-response substrate law.
    \item The later theorem proving that the full displacement-response substrate law collapses benchmark-facingly to the smaller current-body-realizing projection-operator core \(\ProjectionOperatorCore\).
    \item The current theorem proving that the full projection-operator-core body presentation collapses benchmark-facingly to the smaller canonical current-body image core \(\CurrentBodyImageCore\) because off-image kernel directions are benchmark neutral once the canonical image body is fixed.
\end{enumerate}
\end{definition}

\begin{remark}
The present note does not replace those manuscripts. It records the routing consequence of reading them together under the active safeguard against recursive depth drift.
\end{remark}

\section{Why the Live Burden Now Sits Below the Current-Body-Image-Core Frontier}

\begin{proposition}[The live burden is renewed below current-body-image-core restatement]
\label{prop:renewed}
At the current stage of the primitive-reservoir observer-reduction line, the current-body-image-core theorem renews the live benchmark-facing burden below any mere restatement of the canonical current-body image core \(\CurrentBodyImageCore\).
\end{proposition}

\begin{proof}
The new theorem changes benchmark-facing status because it removes the full substrate-body presentation inside the projection-operator core as a live unexplained stopping point. Once the benchmark-facing body content factors through \(\CurrentBodyImageCore\), another manuscript that merely preserves that same smaller image core no longer removes any remaining benchmark-facing freedom. The live burden therefore no longer sits on restating the same image core. It sits on deriving or sharpening that already-active image core from still more primitive explicit substrate structure.
\end{proof}

\begin{definition}[Benchmark-neutral same-image-core relay below the current-body-image-core frontier]
A \emph{benchmark-neutral same-image-core relay below the current-body-image-core frontier} is a manuscript that preserves the same benchmark one-metric output, the same ordinary matter route, the same canonical current-body image core \(\CurrentBodyImageCore\) up to image-core-faithful equivalence, the same current observer-relevant pullback body-response core \(\PullbackBodyResponseCore\), and the same downstream benchmark-facing consequences while merely relocating the unexplained substrate layer deeper, restoring the same full substrate body with benchmark-neutral off-image kernel directions, renaming the same smaller image core, or re-expanding back to the full projection-operator core, the full displacement-response substrate law, or the larger pullback-body-operator, pullback-operator, pullback-quadratic, hidden-response, residual-normal-form, or pre-proto-section presentations without a new derivation, new forcing theorem, or comparably sharp new gain.
\end{definition}

\section{The Current Post-Image-Core Criterion}

\begin{definition}[Admissible next benchmark-facing manuscript below the current-body-image-core frontier]
Below the current current-body-image-core frontier, a manuscript counts as an admissible next benchmark-facing move only if it does at least one of the following:
\begin{enumerate}
    \item derives or justifies the canonical current-body image core \(\CurrentBodyImageCore\) itself from still more primitive explicit substrate structure in a genuinely benchmark-facing way;
    \item derives or justifies some nontrivial constituent of \(\CurrentBodyImageCore\), such as the canonical image body, the restricted projection/canonical lift relation, the exact-shell image, or the response operator restricted to the canonical image, from still more primitive explicit substrate structure in a way that sharpens the present route;
    \item proves a sharper necessity, uniqueness, or compatibility theorem for some nontrivial constituent or relation inside \(\CurrentBodyImageCore\);
    \item does one of the previous three while preserving the same benchmark one-metric package, the same ordinary matter route, and the same claim boundary.
\end{enumerate}
\end{definition}

\begin{theorem}[Next honest benchmark-facing task below the current-body-image-core frontier]
The next honest benchmark-facing manuscript below the current current-body-image-core frontier must derive or justify the smaller image core \(\CurrentBodyImageCore\), or some nontrivial constituent or relation inside it, from still more primitive explicit substrate structure in a benchmark-relevant way, or else prove a genuinely sharper theorem inside that same core. A manuscript that only repeats the same current image-core datum through another deeper relay, restores the already-screened full substrate-body presentation, re-expands the discussion back to the full projection-operator core, the full displacement-response substrate law, or the larger pullback-body-operator, pullback-operator, pullback-quadratic, hidden-response, residual-normal-form, or pre-proto-section presentations without such a new gain, or invokes a side branch without doing the same benchmark-facing work does not satisfy the criterion.
\end{theorem}

\begin{proof}
By Proposition \ref{prop:renewed}, the live burden already lies below any mere restatement of the same smaller image core \(\CurrentBodyImageCore\). Therefore a subsequent manuscript must improve on that already-active core rather than merely accompany it. A deeper-origin manuscript can do so only if it changes benchmark-facing status by more than depth alone. If it simply reproduces the same image core and its same downstream outputs, the routing safeguard classifies it as benchmark neutral. The remaining honest alternatives are therefore exactly those stated in the definition: derive or justify the image core itself, derive or justify some nontrivial constituent or relation inside it, or prove a sharper internal theorem there.
\end{proof}

\section{What the Next Move Should Not Be}

\begin{proposition}[Screened next moves below the current-body-image-core frontier]
Below the current current-body-image-core frontier, none of the following is the default next benchmark-facing move:
\begin{enumerate}
    \item another same-output deeper-origin relay that merely preserves the current image core \(\CurrentBodyImageCore\) up to image-core-faithful equivalence and therefore merely reproduces the same current body-response core \(\PullbackBodyResponseCore\) together with the same downstream benchmark-facing consequences;
    \item another restatement of the full projection-operator core as if the separately listed full substrate-body presentation were again independently live;
    \item another restatement of the full displacement-response substrate law as if the already-forced minimal lift and the already-screened off-image body directions were again primitive stopping points;
    \item another re-expansion back to the larger pullback-body-operator, pullback-operator, pullback-quadratic, hidden-response, residual-normal-form, or pre-proto-section presentations as if those already-spent larger presentations were again independently live burdens;
    \item a side branch that does not derive or justify the same image core, derive or justify a nontrivial constituent or relation inside it, or close a benchmark derivation gate in a genuinely new way;
    \item a manuscript that introduces a second operative metric, enlarges the ordinary matter law, or uses stronger promotion language.
\end{enumerate}
\end{proposition}

\begin{proof}
Item (1) fails the preceding criterion because it leaves the renewed live burden unchanged. Item (2) likewise fails because it merely restores a larger body presentation whose off-image directions the current theorem already classifies as benchmark neutral. Item (3) fails because it reintroduces larger already-spent law data rather than deriving or sharpening the already-active smaller image core. Item (4) fails because it re-expands to larger already-spent presentations rather than improving the live burden. Item (5) does not directly address the live burden at current scope unless it performs the same benchmark-facing work named above. Item (6) violates the fixed claim boundary of the benchmark package and therefore cannot count as a disciplined next step on the active line.
\end{proof}

\begin{remark}
This screening statement is not a denial that those deeper or alternative manuscripts may matter strategically. It is a routing statement: they are not the default way to spend the next benchmark-facing move below the current-body-image-core frontier unless they do the same benchmark-facing work named above.
\end{remark}

\section{Conclusion}

The positive result of this note is routing discipline below the current-body-image-core frontier. The present theorem collapsing the full projection-operator-core body presentation to the smaller canonical current-body image core \(\CurrentBodyImageCore\) is now the live benchmark-facing gain because it removes the larger full substrate-body presentation as a live unexplained stopping point on the recorded route.

The scope remains narrow and honest. The benchmark package is unchanged: the one operative metric is still exactly \(\CompletedMetric\), the ordinary matter route is unchanged, there is no second operative metric, and the low-energy matter law is not enlarged. Nothing here upgrades adoption to a completed final derivation or licenses stronger promotion language.

The next open burden is therefore clear. The next benchmark-facing manuscript should derive or justify the current-body image core itself, or some nontrivial constituent or relation inside it, from still more primitive explicit substrate structure in a way that changes benchmark-facing status. Another same-image-core deeper relay, another restoration of the full substrate-body presentation, a replay of the full projection-operator core, a replay of the full displacement-response substrate law, or a re-expansion back to the larger already-spent presentations is not the clean next manuscript target at current scope.

\end{document}
